\documentclass{article}

\usepackage{tabularx}
\usepackage{booktabs}
\usepackage{longtable}
\usepackage{array}
\usepackage{enumitem}
\usepackage[margin=1in]{geometry}

\title{Specialist Interview Report\\\progname}

\author{\authname}

\date{}

\input{../../Comments.text}
\input{../../Common.text}

% Links into the transcript appendix, as in the SRS.
\newcommand{\intlink}[3]{\hyperref[int:#1-#2]{#3}}
\newcommand{\intref}[2]{\intlink{#1}{#2}{#1, #2}}

\begin{document}

\maketitle

\begin{table}[hp]
\caption{Revision History} \label{TblRevisionHistory}
\begin{tabularx}{\textwidth}{llX}
\toprule
\textbf{Date} & \textbf{Developer(s)} & \textbf{Change}\\
\midrule
11 November 2026 & Rayan Nasrallah & First version: methodology, questions,
subjects, empathy map analysis, results and limitations, with the transcript
as an appendix (\#262).\\
\bottomrule
\end{tabularx}
\end{table}

\newpage
\tableofcontents
\newpage

\section{Purpose}

This report is the fall custom document for \progname{}. It describes the
interview we held with three domain specialists on October 1, 2026, and
analyses what they told us with an empathy map. The Software Requirements
Specification (SRS) already traces each interview statement to the
requirements it shaped (SRS Appendix A, SRC-005). This report does not repeat
that trace. Instead, it looks at the interview as a whole: what the
specialists said, thought, did and felt, which themes run across their
answers, and what we still need to ask them. Those open questions feed the
next session, where the same specialists will try the proof of concept.

\section{Methodology}

\subsection{Interview type}

We ran one semi-structured group interview. A semi-structured interview gives
a fixed core of questions, so every topic we needed gets raised, while
follow-up probes let the specialists take the discussion where their
experience points. A structured interview would have lost the reasoning
behind each answer, and an unstructured one would have been hard to keep on
track with three people in 45 minutes.

\subsection{Procedure}

\begin{itemize}
\item \textbf{Format:} a remote video call of 43 minutes, recorded and
  transcribed automatically.
\item \textbf{Interviewers:} Rayan Nasrallah (facilitator) and Moustafa
  Moustafa.
\item \textbf{Structure:} a 10-minute presentation of the planned application
  (guided lessons, creative drawing space, practice history for adults, and
  two possible future features), then the questions in
  Section~\ref{sec:questions}.
\item \textbf{Consent:} each specialist signed a consent form before the
  session, and all three agreed during the session to the use of the
  transcript for eliciting requirements (\intlink{Rayan}{41:51}{41:51--42:48}).
\item \textbf{Follow-up:} the questions we ran out of time for went into a
  short written questionnaire, answered by two of the three specialists on
  October 2 and 4, 2026 (SRS Appendix, SRC-009).
\end{itemize}

\subsection{Data preparation}

The automatic transcript was de-identified at the specialists' request:
names, institutions, locations and a personal health detail were removed or
replaced with bracketed placeholders. Mis-heard words were corrected and
filler words removed. Each turn keeps its timestamp from the recording, so
every claim in this report can be checked against the transcript in
Appendix~\ref{app:interview}.

\subsection{Analysis method: empathy map}

We analysed the transcript with an empathy map, a design thinking tool that
sorts what is known about a group of people into four quadrants
\cite{gray2010gamestorming, gibbons2018empathy}:

\begin{itemize}
\item \textbf{Says:} what they stated directly.
\item \textbf{Thinks:} the beliefs and priorities behind what they said.
\item \textbf{Does:} what they reported doing in their own work.
\item \textbf{Feels:} the reactions and concerns they expressed.
\end{itemize}

We chose an empathy map for two reasons. First, the SRS already links single
statements to requirements, so a second statement-by-statement coding would
add nothing. An empathy map instead shows the specialists' point of view as a
whole, which is what we need to plan the next session with them. Second, it
separates what was said from what was reported as practice, which matters
because the specialists' experience comes from their own projects, not from
observing our users.

The steps were:
\begin{enumerate}
\item Read the full transcript and mark every turn by a specialist that
  carries content (not greetings or audio checks).
\item Place each marked turn in one or more quadrants, keeping its
  timestamp. ``Says'' holds paraphrases of direct statements. ``Thinks'' and
  ``Feels'' only hold beliefs and reactions the specialists put into words; we
  did not guess at unspoken ones.
\item Group the entries into themes, then lay the map out with one row per
  theme, so each statement sits beside the belief, practice and reaction
  behind it.
\item For each theme, note where the specialists disagreed, and what we still
  need to ask.
\end{enumerate}

We built one map for the group rather than one per specialist; the empathy
map section explains why.

\section{Subjects}
\label{sec:subjects}

Three domain specialists took part. They research and develop software for
children with autism, ADHD and other learning needs, and hold doctorates in
human-computer interaction and related engineering fields. They work at the
same research centre outside Canada. They preferred not to be named, so this
report refers to them as P1, P2 and P3 and gives no further detail that could
identify them. With only three people from one centre, we report their
background as a group rather than per person.

\begin{table}[h]
\centering
\begin{tabularx}{\textwidth}{lX}
\toprule
\textbf{Code} & \textbf{Role in the session} \\
\midrule
P1 & Specialist. Asked the first questions after the presentation and led the
answers on progression, child independence and input. \\
P2 & Specialist. Led the answers on population, accessibility guidelines and
audio feedback. \\
P3 & Specialist. Raised the question of pen or finger input. \\
\bottomrule
\end{tabularx}
\caption{Subjects}
\end{table}

\section{Questions}
\label{sec:questions}

\subsection{Questions asked in the session}

We had prepared a guide of 16 main questions. After the presentation, the
specialists raised their own questions first, and the discussion followed
them. Table~\ref{tbl:asked} lists the questions as they were actually asked,
by us or by the specialists, in order.

\begin{longtable}{p{0.08\textwidth}p{0.15\textwidth}p{0.68\textwidth}}
\toprule
\textbf{Time} & \textbf{Asked by} & \textbf{Question} \\
\midrule
\endhead
\intlink{P1}{12:07}{12:07} & P1 & Did you co-design the app with educators,
therapists or other end users? \\
\intlink{P1}{13:42}{13:42} & P1 & (Comment) Camera finger tracing may target a
different skill. \\
\intlink{Moustafa}{14:55}{14:55} & Moustafa & Is camera finger tracing helpful
if we do implement it? \\
\intlink{Rayan}{15:48}{15:48} & Rayan & With limited time, should we focus on
shapes and tracing before the camera feature? \\
\intlink{P2}{16:45}{16:45} & P2 & Are you targeting specific groups of
children, or a broader population? \\
\intlink{Moustafa}{20:14}{20:14} & Moustafa & So the app should not be limited
to children with autism or ADHD? \\
\intlink{Rayan}{22:07}{22:07} & Rayan & Would you take the app further in your
own research, or help observe children using it? \\
\intlink{P3}{24:34}{24:34} & P3 & Will children draw with a pen or with their
finger? \\
\intlink{Rayan}{27:31}{27:31} & Rayan & Should the app decide when a child is
ready to move to the next level, or leave that to the supervising adult? What
determines readiness? \\
\intlink{Rayan}{31:10}{31:10} & Rayan & What progress information would help a
parent or caregiver judge whether the child is ready to move on? \\
\intlink{Rayan}{32:12}{32:12} & Rayan & Can you connect us with occupational
therapists or educational specialists? \\
\intlink{Moustafa}{33:34}{33:34} & Moustafa & Should a child be able to use the
app alone, or should we assume an adult is present? How should the app give
feedback and instructions? \\
\intlink{Rayan}{37:18}{37:18} & Rayan & Is spoken audio feedback a good
idea? \\
\intlink{Rayan}{38:58}{38:58} & Rayan & Which design choices cause problems for
autistic children or children with ADHD, and which practices should we
follow? \\
\intlink{Rayan}{40:44}{40:44} & Rayan & Any other comments or patterns we
should know about? \\
\bottomrule
\caption{Questions asked in the session}
\label{tbl:asked}
\end{longtable}

\subsection{Coverage of the prepared guide}

Table~\ref{tbl:coverage} shows which topics of the prepared guide were covered
in the session, which went to the follow-up questionnaire, and which were not
covered.

\begin{table}[h]
\centering
\begin{tabularx}{\textwidth}{Xl}
\toprule
\textbf{Guide topic} & \textbf{Covered by} \\
\midrule
First reactions & Session \\
Whether and how children practise these skills now & Not covered \\
Existing apps and tools & Not covered \\
What adults find hardest when guiding a child & Not covered \\
School tool, home tool or both & Partly (P2, 35:36) \\
Adult present or child alone & Session \\
Devices and input methods & Session (pen or finger) \\
What motivates or frustrates children & Questionnaire \\
How to give feedback & Session (audio) and questionnaire \\
How adults decide a child is ready for less support & Session \\
Design choices that cause problems & Session \\
Whether cleaning up free drawings helps & Questionnaire \\
Useful practice information for adults & Session \\
Camera finger tracing & Session \\
Letter tracing & Questionnaire \\
Top priorities & Not covered \\
\bottomrule
\end{tabularx}
\caption{Coverage of the prepared interview guide}
\label{tbl:coverage}
\end{table}

\section{Raw Data}

The de-identified transcript is in Appendix~\ref{app:interview}.

\section{Analysis}

\subsection{Empathy map}

Table~\ref{tbl:map} is the empathy map. Empathy maps are often built for
one user at a time and then combined into one map for a user type
\cite{gibbons2018empathy}. Instead, we built a single map for all three
specialists, for two reasons. First, they answered as a group: they built on each other's
answers, so splitting the session by person would separate points that were
made together. Second, the specialists did not speak equally. P3 spoke on one
topic (pen or finger input), so a map for P3 alone would have about two
entries. Keeping one map, with every entry labelled by speaker, shows who said
what without three thin maps. Section~\ref{sec:themes} shows who led each
theme.

A classic empathy map has four quadrants. Here each row is one theme, so a statement sits beside the belief
behind it, what the specialists do in their own work, and how they felt about
it. Each entry is a short paraphrase with the speaker and a link to the
transcript. A dash marks a quadrant the interview gave no evidence for; those
gaps are discussed in Section~\ref{sec:results}.

\begingroup
\small
\renewcommand{\arraystretch}{1.2}
\begin{longtable}{*{4}{>{\raggedright\arraybackslash}p{0.22\textwidth}}}
\toprule
\textbf{Says} & \textbf{Thinks} & \textbf{Does} & \textbf{Feels} \\
\midrule
\endhead
\bottomrule
\endfoot
\multicolumn{4}{l}{\textbf{Theme 1: Adults decide, the app supports}} \\*
An adult who knows the child should choose the level
(\intref{P1}{28:37}). Let the teacher or caregiver change the activity and
see progress (\intref{P2}{29:18}).
& Having the app decide means assessment, and a progress view raises the
question of what assessment it applies (\intref{P1}{28:37};
\intref{P2}{29:18}).
& Built a vocabulary app where the teacher personalized each child's path
through a separate adult view (\intref{P2}{29:18}).
& Cautious: automatic progression is ``more advanced'' and too complex for
now (\intref{P1}{28:37}). \\
\midrule
\multicolumn{4}{l}{\textbf{Theme 2: The child should not need an adult in the room}} \\*
A child who cannot read should be able to enter a lesson, practise, restart
and move around using images and icons (\intref{P1}{34:13}).
& Children are used to icon-based games; adults handle setup
(\intref{P1}{34:13}). The child-and-caregiver workflow is the team's choice
(\intref{P2}{35:36}).
& --
& Concern that parents get annoyed if they must keep stepping in
(\intref{P1}{34:13}). \\
\midrule
\multicolumn{4}{l}{\textbf{Theme 3: Design the details with the people who know the children}} \\*
Which progress measures help should be decided with occupational therapists
or educators (\intref{P2}{31:42}). Use spoken audio that can be turned off
(\intref{P2}{37:39}).
& Educators and therapists, not designers, know the children day to day
(\intref{P1}{36:26}). Offering more than one modality is good inclusive
design (\intref{P2}{37:39}).
& Design praise wording and colours from the literature and therapists'
feedback, bit by bit (\intref{P1}{36:26}). Gave their app voice input that
users can turn off (\intref{P2}{37:39}).
& -- \\
\midrule
\multicolumn{4}{l}{\textbf{Theme 4: Accessible and broad from the start}} \\*
Start with WCAG, then autism design heuristics (\intref{P2}{39:26}). Don't
limit the app to one diagnosis (\intref{P1}{20:27}).
& Designing for disability serves everyone; the skills apply to all children,
at different stages (\intref{P2}{19:05}; \intref{P2}{20:47}).
& Teaches accessible design using WCAG (\intref{P2}{39:26}).
& P2 related to the problem from their own difficulty drawing lines as a
child (\intref{P2}{16:45}). \\
\midrule
\multicolumn{4}{l}{\textbf{Theme 5: The input tool shapes the skill}} \\*
A finger is easiest (\intref{P3}{25:36}). Fine motor practice needs a held
tool, from large to small (\intref{P1}{26:46}).
& How a child holds a pen affects their drawing and the app's usability
(\intref{P3}{25:36}).
& --
& -- \\
\midrule
\multicolumn{4}{l}{\textbf{Theme 6: Core first, with care around children}} \\*
Focus on touchscreen tracing and add camera tracing later
(\intref{P1}{16:13}; \intref{P2}{16:45}). Testing with children needs ethics
approval; meeting experts does not (\intref{P2}{22:56}; \intref{P1}{33:00}).
& Camera tracing trains gross rather than fine motor skills, but can help a
child visualize the shape (\intref{P1}{13:42}; \intref{P1}{15:12}).
& Obtains ethics approval for studies with children, which takes about three
to four months (\intref{P2}{22:56}).
& Willing to keep helping: to meet again, guide a pilot and share contacts
(\intref{P1}{24:03}; \intref{P1}{41:08}; \intref{P2}{22:56}). \\
\midrule
\multicolumn{4}{l}{\textbf{Overall}} \\*
-- & -- & --
& Positive about the project: ``very nice'', ``inspiring'', ``much needed'',
``timely'' (\intref{P2}{16:45}; \intref{P1}{24:03}; \intref{P1}{41:08}). \\
\caption{Empathy map of the three specialists, by theme}
\label{tbl:map}
\end{longtable}
\endgroup

\subsection{Themes}
\label{sec:themes}

\begin{enumerate}
\item \textbf{Adults decide, the app supports.} Every quadrant agrees: the
  adult chooses the level, because anything else becomes an assessment.
\item \textbf{The child should not need an adult in the room.} This pulls
  against theme 1: the adult controls progression but must not have to sit
  through every session.
\item \textbf{Design the details with the people who know the children.}
  Feedback wording, colours, audio and progress measures should come from
  occupational therapists and educators.
\item \textbf{Accessible and broad from the start.} WCAG first, then autism
  heuristics, for all children rather than one diagnosis.
\item \textbf{The input tool shapes the skill.} The one clear disagreement:
  P3 (and later R2) for a finger, P1 for a held tool when fine motor practice
  is the goal.
\item \textbf{Core first, with care around children.} Touchscreen tracing
  before camera tracing, and expert feedback rather than testing with
  children.
\end{enumerate}

Table~\ref{tbl:leads} shows which specialist led each theme.

\begin{table}[h]
\centering
\begin{tabularx}{\textwidth}{lX}
\toprule
\textbf{Theme} & \textbf{Who led it} \\
\midrule
1. Adults decide & P1 first, then P2 from a past project \\
2. Child without an adult & P1; P2 added that the workflow is the team's choice \\
3. Design with those who know the children & P1 (feedback wording and colours) and P2 (audio, progress measures) \\
4. Accessible and broad & P2; P1 agreed \\
5. Input tool & P3 raised it; P1 answered \\
6. Core first, care around children & P1 (camera tracing, ethics); P2 agreed and gave the ethics timeline \\
\bottomrule
\end{tabularx}
\caption{Who led each theme}
\label{tbl:leads}
\end{table}

\subsection{How the themes were used}

We built what the specialists told us, as our supervisor advised. Themes 1, 2,
4 and 6 became requirements and constraints in the SRS. For theme 5 we
followed P1: the stylus is the main input because fine motor practice is the
goal, and a finger still works. Theme 3 is only partly met, and here we
consulted our supervisor. The specialists suggested deciding progress measures
with occupational therapists (\intref{P2}{31:42}). Our supervisor advised
against waiting on one for the values that depend on the child, such as how
accurate a drawing must be to pass a stage and how many failed attempts end a
lesson, because these differ from child to child, so an occupational therapist
could not give one right value either (supervisor meeting, October 14, 2026).
Instead, the adult who knows the child sets these values for that child, within
a minimum and maximum the team defines and shows with example drawings, and
the midpoint is the default. This avoids the risk of the team fixing values
that suit some children but not others, which we are not qualified to judge.
It also follows theme 1, since the adult already decides the child's level. In
the same way, the adult chooses which recorded values their progress view
shows. Feedback wording and colours will be reviewed with the specialists in
the follow-up session. The SRS records the trace from each
statement to its requirements (SRS Appendix A, SRC-005.01 to SRC-005.13).

\section{Results and Next Steps}
\label{sec:results}

The map shows that the specialists' view is consistent: a child-friendly app
that the child can use alone, controlled by the adult who knows the child,
accessible from the start, and designed with the people who work with
children every day. Their main reservation is any feature that looks like an
assessment.

The dashes in the map show what we did not learn. The ``Does'' column is empty
for child independence and input, because the guide's questions on how
children practise now and which tools exist were not reached. The ``Feels''
column says little about feedback and input, and nothing about how children
react, because the specialists spoke from their own projects.

We will return to the same specialists with the proof of concept, a working
lesson they can try on their own time. That session will check our reading of
each theme, and ask about the gaps:
\begin{itemize}
\item Can a child who cannot read get through the lesson with the icons and
  audio we chose (themes 1 and 2)?
\item Are the feedback wording, colours and sounds suitable (theme 3)?
\item With the stylus and a finger both available, does the stylus still seem
  the right main input (theme 5)?
\item Does the adult's progress view read as an assessment (themes 1 and 3)?
\item How do children practise these skills now, and which existing tools fall
  short (not covered in the first session)?
\end{itemize}
In that session we will run the demo first and explain our decisions only
after it ends, so the specialists react to the app rather than to our
reasoning.

\section{Limitations}

\begin{itemize}
\item \textbf{Specialists only.} We interviewed three specialists. We did not
  interview parents, teachers or occupational therapists directly, and we did
  not involve children.
\item \textbf{Expert opinion, not observation.} The specialists described
  their own experience and projects. Nothing here comes from observing
  children use \progname{}.
\item \textbf{One session, one centre.} All three work at the same centre, in
  one 43-minute session, so their views may share a common background.
\item \textbf{Group and presentation effects.} In a group, one answer can
  shape the next, and presenting our design first may have led the
  specialists to react to our features rather than raise needs we missed.
\item \textbf{Inferred quadrants.} ``Thinks'' and ``Feels'' rely on what the
  specialists put into words. We did not infer unspoken feelings, so these
  quadrants are incomplete.
\item \textbf{Coverage.} Four topics of the guide were not covered (see
  Table~\ref{tbl:coverage}), and only two of three specialists answered the
  follow-up questionnaire.
\end{itemize}

\begin{thebibliography}{9}
\bibitem{gray2010gamestorming} D.~Gray, S.~Brown and J.~Macanufo,
  \emph{Gamestorming: A Playbook for Innovators, Rulebreakers, and
  Changemakers}. Sebastopol, CA: O'Reilly Media, 2010.
\bibitem{gibbons2018empathy} S.~Gibbons, ``Empathy mapping: The first step in
  design thinking,'' Nielsen Norman Group, Jan.~2018. [Online]. Available:
  \url{https://www.nngroup.com/articles/empathy-mapping/}
\end{thebibliography}

\newpage
\appendix

\section{Generative AI Use Disclosure}
\label{app:ai-use}

This appendix describes how generative AI tools assisted in preparing and
reviewing this document. The decisions in this report came from the specialist
interview, which Moustafa and I conducted, and from our meetings with our
supervisor and teaching assistant, including the choice of an empathy map as
the analysis method. AI tools helped carry out and express that work; they did
not make decisions for us. The author remains responsible for the content of
the final document.

\subsection{AI Tools Used}

\begin{itemize}
  \item \textbf{Anthropic's Claude Opus 5.5}, accessed through Claude Code at
  low effort.
\end{itemize}

\subsection{How and Where Used}

The uses below materially contributed to this document.

\begin{itemize}
  \item \textbf{Organizing the interview data.} AI helped sort through the
  interview transcript and locate the specific statements, with their
  timestamps, behind each empathy map entry, each question in
  Table~\ref{tbl:asked} and each theme. This affected Sections 4 and 6.

  \item \textbf{Drafting and wording.} AI turned my notes and decisions into
  draft wording for each section, and helped improve sentence structure,
  clarity and flow. I then directed changes, such as laying the empathy map out
  with one row per theme, explaining why we built one map for all three
  specialists, and adding the supervisor's advice in Section 6.3.

  \item \textbf{Formatting.} AI helped with the \LaTeX{} formatting, including
  the tables and the links into the transcript.

  \item \textbf{Repository management.} Through Claude Code, AI helped prepare
  the branch and commit for this document.
\end{itemize}

\subsection{Verification of AI-Generated Content}

Every AI-assisted passage was reviewed against my intended meaning and the
decisions the team had agreed on before it was included. Any wording that
introduced claims the specialists did not make, or commitments the team had
not agreed on, was removed or rewritten.

Each empathy map entry and each listed question was checked against the
original transcript to confirm that the cited statement and timestamp say what
this report claims. The ``Thinks'' and ``Feels'' columns only hold beliefs and
reactions the specialists put into words. The references were checked against
their sources, and the report was checked for details that could identify the
specialists.

\textbf{The author is responsible for the accuracy, completeness, consistency,
and appropriateness of this document, regardless of whether material was
written by a team member or with the help of an AI system.}

\newpage
\input{../../SRS/InterviewTranscript.text}

\end{document}
